% Generated draft: requires mathematical and source review.
\documentclass[11pt]{article}
\usepackage[margin=25mm]{geometry}
\usepackage{fontspec}
\setmainfont{Latin Modern Roman}
\usepackage{amsmath,amssymb,mathtools,unicode-math}
\setmathfont{Latin Modern Math}
\usepackage{enumitem,needspace,xurl,hyperref}
\setlength{\parindent}{0pt}
\setlength{\parskip}{5pt}
\setlength{\emergencystretch}{3em}
\newcommand{\N}{\mathbb{N}}
\newcommand{\Z}{\mathbb{Z}}
\newcommand{\Q}{\mathbb{Q}}
\newcommand{\R}{\mathbb{R}}
\newcommand{\C}{\mathbb{C}}
\newcommand{\F}{\mathbb{F}}
\newcommand{\A}{\mathbb{A}}
\newcommand{\PP}{\mathbb{P}}
\newcommand{\E}{\mathbb{E}}
\newcommand{\eps}{\varepsilon}
\newcommand{\dd}{\,\mathrm{d}}
\begin{document}
\section*{UnsolvedMath ID: 30005392}
\textbf{Source problem code:} catalogue entry
\subsection{Definitions and mathematical statement}
Let \(p\) be a prime, and let \(\mathcal A_g\) denote the moduli space of principally polarised abelian varieties of dimension \(g\) in characteristic \(p\). The source gives no restriction on \(p\) or explicit range for \(g\). Write \(x=(A_x,\lambda_x)\), where \(\lambda_x\) is a principal polarisation. For an abelian variety \(A\), let \(\nu(A)\) denote the Newton polygon of its \(p\)-divisible group \(A[p^\infty]\). The Newton stratum containing \(x\) is the locus
\[
\mathcal N_x=\{y\in\mathcal A_g:\nu(A_y)=\nu(A_x)\}.
\]
Here the Newton polygon at a point is computed after passage to an algebraic closure of its residue field.

The supplied formulation takes \(x\in\mathcal A_g(\F_p)\) and describes its Hecke orbit as
\[
\mathcal H_{\mathrm{src}}(x)=\bigl\{y=(A_y,\lambda_y)\in\mathcal A_g(\F_p):
(A_y,\lambda_y)\text{ is related to }(A_x,\lambda_x)
\text{ by symplectic isogenies}\bigr\}.
\]
The recorded conjectural assertion is that, for every such \(x\),
\[
\overline{\mathcal H_{\mathrm{src}}(x)}^{\,\mathcal N_x}=\mathcal N_x,
\]
where the bar denotes Zariski closure in the Newton stratum.

\textbf{Unresolved qualifications.} This is a transcription of the target assertion, not a fully specified formulation. The supplied text repeatedly writes \(\F_p\), without an algebraic closure. Taken literally, its orbit is contained in the finite set \(\mathcal A_g(\F_p)\); such a set cannot be dense in a positive-dimensional Newton stratum. It must therefore be checked whether geometric points over \(\overline{\F}_p\) were intended, and over which field the isogenies are allowed to be defined. No replacement of \(\F_p\) by \(\overline{\F}_p\) has been silently made here.

The expression ``symplectic isogenies'' also lacks a convention in the supplied record. In particular, the record does not specify whether the isogenies must have degree prime to \(p\), whether polarisation similitudes are allowed, or what multipliers are permitted. For an isogeny \(f:A_x\to A_y\), a polarisation-similitude condition would involve \(f^\vee\lambda_y f=c\lambda_x\), but neither this convention nor the allowed values of \(c\) are fixed by the record. Consequently, \(\mathcal H_{\mathrm{src}}(x)\) retains the source's unresolved terminology. The Newton-stratum target is preserved; replacing it by another locus would require source verification.
\subsection{Short English statement}
The recorded conjecture asks whether the orbit of a principally polarised abelian variety under the specified symplectic isogenies is Zariski dense among varieties with the same Newton polygon. The field of points and the precise isogeny convention remain unresolved. The supplied status and verification labels are curation claims, not independent verification of this draft.
\subsection{Sources}
\begin{itemize}
\item Local UnsolvedMath registry, permanent ID \(30005392\), supplied complete source record, title ``Oort's Hecke-Orbit Conjecture'', source problem code \texttt{catalogue entry}.
\item \emph{Arithmetic of Shimura Varieties} (2023), Oberwolfach Reports, DOI 10.4171/owr/2023/5, pp.\ 261--326. Supplied source URL and literature-source entry ``Oberwolfach Report 5/2023'': \url{https://doi.org/10.4171/owr/2023/5}. The supplied excerpt attributes the presentation ``Hecke orbits on Shimura varieties of Hodge type'' to Pol van Hoften, joint work with Marco D'Addezio, and identifies the conjecture as Problem 15 on Oort's 1995 list of open problems in algebraic geometry, referenced there as [14]. The bibliographic entry for [14] is not supplied.
\item D'Addezio and van Hoften, \emph{Hecke orbits on Hodge-type Shimura varieties}: \url{https://arxiv.org/abs/2205.10344}. This is a supplied literature reference; its contents have not been independently checked for this draft.
\item Supplied literature assessment dated \texttt{2026-08-22}: status ``partially solved''. The assessment states that Oort's Hecke orbit conjecture is proved for Siegel varieties and for Hodge-type Shimura data under a lower bound on \(p\), while the unrestricted conjecture in the extracted statement is not settled. This is supplied dated curation, not independently verified; its scope and compatibility with the extracted \(\mathcal A_g\) statement require review. The background explicitly requests independent verification of time-sensitive status and cited claims.
\item The registry labels the statement \texttt{corrected\_verified} and reports that it was reconstructed and checked against the official Oberwolfach report. These are supplied provenance claims, not verification performed in preparing this statement-only draft.
\end{itemize}
\end{document}
